% Discussion chapter

\section{Interpretation of Results}
\subsection{Model Performance}
Interpret what the performance metrics indicate about the feasibility of ML-based shadow fleet detection.

\subsection{Feature Insights}
Discuss what the important features reveal about shadow fleet vessel behavior.

\section{Comparison with Related Work}
Compare your results with similar research in maritime anomaly detection.

\section{Practical Implications}
\subsection{Operational Deployment}
Discuss how this system could be deployed in real-world maritime monitoring.

\subsection{Benefits and Use Cases}
Potential applications for maritime authorities, port operators, etc.

\subsection{Integration with Existing Systems}
How this could complement existing maritime surveillance infrastructure.

\section{Challenges and Limitations}
\subsection{Data Quality Issues}
Discuss limitations related to AIS data quality, spoofing, and gaps.

\subsection{Evolving Tactics}
How shadow fleet operators may adapt their tactics over time.

\subsection{False Positive Management}
Challenges in minimizing false alarms while maintaining detection sensitivity.

\subsection{Legal and Ethical Considerations}
Privacy concerns, due process, and the appropriate use of automated detection.

\section{Generalizability}
\subsection{Applicability to Other Regions}
Could this approach work in other maritime regions?

\subsection{Transferability to Other Vessel Types}
Could similar methods detect other types of illegal maritime activities?

\section{Lessons Learned}
Key insights gained during the research process.
